\chapter{PENDAHULUAN}

\section{Latar Belakang}
Emas merupakan salah satu aset yang banyak diperdagangkan di pasar keuangan global dan sering digunakan sebagai instrumen lindung nilai (\emph{hedge}) ketika ketidakpastian ekonomi meningkat. Dalam pasar valuta asing, emas diperdagangkan melalui XAUUSD, yaitu harga satu troy ounce emas dalam dolar Amerika Serikat. Pergerakan harganya berkaitan dengan berbagai kondisi ekonomi dan pasar. \textcite{dinh2022economic} menunjukkan hubungan volatilitas dan korelasi logam mulia dengan pasar saham, neraca perdagangan, kepercayaan konsumen, jumlah uang beredar, dan inflasi. \textcite{wang2023asymmetric} juga menemukan hubungan asimetris antara imbal hasil logam mulia, harga minyak mentah, dan nilai tukar dolar Amerika Serikat dengan mempertimbangkan volatilitas pasar serta suku bunga. Pada masa guncangan global, hubungan lintas pasar tersebut dapat menguat, sebagaimana ditunjukkan oleh \textcite{yousaf2022information} melalui transmisi informasi antara emas dan minyak mentah selama wabah COVID-19. Banyaknya faktor yang terlibat membuat perkiraan pergerakan harga emas menjadi persoalan yang penting sekaligus menantang.

Untuk berbagai keperluan analisis, informasi mengenai arah pergerakan harga berikutnya dapat lebih berguna daripada perkiraan nilai harga secara tepat. Oleh karena itu, sejumlah penelitian merumuskan prediksi harga sebagai masalah klasifikasi arah. \textcite{sadorsky2021predicting} memprediksi arah harga emas dan perak dengan model logit dan beberapa pengklasifikasi berbasis pohon menggunakan indikator teknikal. Model berbasis pohon menghasilkan kinerja yang lebih baik daripada model linear pada sebagian besar konfigurasi yang diuji. \textcite{basher2022forecasting} kemudian memasukkan suku bunga, inflasi, dan volatilitas pasar dalam prediksi arah harian emas dan Bitcoin. Kedua penelitian tersebut menunjukkan bahwa jenis model dan kelompok informasi yang digunakan dapat memengaruhi kinerja klasifikasi arah.

Informasi untuk memprediksi arah harga dapat berasal dari riwayat harga maupun kondisi di luar pasar. Indikator teknikal mengubah riwayat harga menjadi representasi momentum atau tren yang lebih ringkas. Sebagai contoh, \textcite{mamoudan2023hybrid} memanfaatkan turunan MACD, \emph{moving average}, dan Ichimoku untuk mengklasifikasikan sinyal teknikal pada pasar emas sebagai \emph{FAKE} atau \emph{REAL}. Di sisi lain, variabel makroekonomi membawa informasi mengenai kondisi ekonomi. \textcite{sobti2021how} menunjukkan bahwa pengumuman makroekonomi memengaruhi proses pembentukan harga emas sepanjang hari perdagangan. Temuan ini mendukung pengujian informasi kalender ekonomi pada resolusi intrahari, dengan tetap memperhatikan waktu ketika informasi tersebut benar-benar tersedia.

Penambahan fitur belum tentu selalu meningkatkan kinerja karena dapat menambah dimensi, derau, dan risiko pemodelan berlebih (\emph{overfitting}). \textcite{taneva2025framework} telah memadukan data finansial, ekonomi, dan sentimen untuk meramalkan harga emas harian, sedangkan \textcite{velez2026predictive} membandingkan 32 kombinasi model dan fitur untuk peramalan harga emas bulanan. Kedua studi tersebut memperlihatkan pentingnya penggabungan dan perbandingan fitur. Namun, dalam batas studi yang ditelaah, perubahan kinerja akibat penambahan kelompok harga, teknikal, dan makroekonomi secara bertahap belum diukur untuk klasifikasi tiga arah XAUUSD H1 pada observasi uji yang sama.

Resolusi satu jam (H1) dipilih untuk mengamati dinamika intrahari dengan jumlah pengamatan yang lebih banyak daripada data harian. Penggunaan data H1 untuk XAUUSD memiliki preseden pada penelitian \textcite{arif2026multivariate}, yang meramalkan harga penutupan dengan LSTM dan validasi deret waktu bergulir. Pada konteks berbeda, \textcite{hisyam2025implementasi} menggunakan XGBoost dan \emph{walk-forward validation} untuk meramalkan harga emas Antam harian. Penelitian ini menerapkan Regresi Logistik Multinomial sebagai pembanding linear dan XGBoost sebagai pembanding nonlinear pada setiap konfigurasi fitur. Kedua model membantu membaca perubahan kinerja setelah penambahan fitur; perbedaan skor antarmodel ditafsirkan secara deskriptif, bukan sebagai uji linearitas data.

Pengukuran perubahan kinerja juga harus mengikuti urutan waktu. Pembagian data deret waktu secara acak dapat menyebabkan kebocoran informasi masa depan (\emph{look-ahead bias}) dan menghasilkan penilaian yang tidak mewakili penggunaan pada data baru. Penelitian ini karena itu menggunakan evaluasi \emph{walk-forward} dengan jeda antara data pelatihan dan pengujian. Macro-F1 digunakan sebagai metrik utama karena distribusi kelas BUY, HOLD, dan SELL dapat tidak seimbang, sedangkan balanced accuracy digunakan sebagai metrik pendamping.

Berdasarkan uraian tersebut, penelitian ini mengukur perubahan kinerja klasifikasi arah XAUUSD H1 setelah fitur teknikal dan makroekonomi ditambahkan secara bertahap. Konfigurasi harga (K1), harga dan teknikal (K2), serta harga, teknikal, dan makroekonomi (K3) dibandingkan menggunakan Regresi Logistik Multinomial dan XGBoost pada periode uji kronologis yang sama. Perubahan $K2-K1$ dan $K3-K2$ dibaca dalam masing-masing model, sedangkan perbedaan respons kedua model dijelaskan secara deskriptif. Kontribusi fitur ditafsirkan secara kondisional terhadap data, model, periode, dan prosedur penyetelan yang digunakan. Konfigurasi fitur makroekonomi (K3) dievaluasi menggunakan arsip kalender beku MetaTrader 5, sedangkan keterbatasan \emph{snapshot} retrospektif tanpa arsip \emph{vintage} independen membatasi interpretasi kinerja untuk penerapan waktu nyata.

\section{Rumusan Masalah}
Berdasarkan latar belakang tersebut, rumusan masalah penelitian ini adalah sebagai berikut.
\begin{enumerate}
    \item Bagaimanakah kinerja klasifikasi arah XAUUSD H1 pada konfigurasi fitur harga (K1), harga dan teknikal (K2), serta harga, teknikal, dan makroekonomi (K3) ketika masing-masing diterapkan pada Regresi Logistik Multinomial dan XGBoost?
    \item Bagaimanakah besar dan konsistensi perubahan kinerja setelah penambahan fitur teknikal dan fitur makroekonomi apabila dievaluasi pada beberapa periode uji berurutan dan pada kedua model yang digunakan?
\end{enumerate}

\section{Tujuan}
Penelitian ini bertujuan untuk:
\begin{enumerate}
    \item mengukur kinerja klasifikasi arah XAUUSD H1 secara \emph{out-of-sample} untuk K1, K2, dan K3 pada Regresi Logistik Multinomial serta XGBoost; dan
    \item mengukur serta menganalisis besar dan konsistensi perubahan kinerja antarkonfigurasi pada beberapa periode uji berurutan menggunakan Regresi Logistik Multinomial dan XGBoost.
\end{enumerate}

\section{Manfaat}
Hasil penelitian ini diharapkan memberikan manfaat sebagai berikut.
\begin{enumerate}
    \item Memberikan bukti terukur bagi bidang komputasi cerdas mengenai perubahan kinerja klasifikasi arah XAUUSD H1 setelah penambahan fitur teknikal dan makroekonomi dalam batas data, model, dan periode pengujian.
    \item Menyediakan prosedur perbandingan konfigurasi fitur bertahap dan evaluasi kronologis yang dapat diadaptasi oleh penelitian selanjutnya.
    \item Memberikan gambaran bagi peneliti dan praktisi mengenai perubahan kinerja yang diperoleh ketika kelompok fitur teknikal dan makroekonomi ditambahkan secara bertahap.
\end{enumerate}

\section{Batasan Masalah}
Agar pembahasan tetap terarah, penelitian ini dibatasi sebagai berikut.
\begin{enumerate}
    \item Objek penelitian adalah XAUUSD dengan timeframe satu jam (H1), sehingga simpulan dibatasi pada instrumen dan timeframe tersebut.
    \item Target pemodelan terdiri atas tiga label arah, yaitu BUY, HOLD, dan SELL, yang dibentuk dari perubahan harga penutupan satu bar berikutnya. Label tersebut tidak dimaksudkan sebagai rekomendasi transaksi.
    \item Model yang digunakan adalah Regresi Logistik Multinomial dan XGBoost.
    \item Fitur makroekonomi berasal dari data terstruktur kalender ekonomi MetaTrader 5, meliputi nilai indikator dan penanda peristiwa. Teks berita dan analisis sentimen tidak termasuk dalam penelitian.
    \item Kinerja diukur menggunakan macro-F1 dan balanced accuracy melalui evaluasi \emph{walk-forward}.
    \item Analisis bersifat prediktif, tidak menyimpulkan hubungan sebab-akibat ekonomi, menjanjikan keuntungan transaksi, atau menghasilkan sistem eksekusi order otomatis.
\end{enumerate}

\section{Sistematika Pembahasan}
Skripsi ini disusun dalam enam bab sebagai berikut.
\begin{enumerate}
    \item BAB 1 PENDAHULUAN menjelaskan latar belakang, rumusan masalah, tujuan, manfaat, batasan masalah, dan sistematika pembahasan.
    \item BAB 2 LANDASAN KEPUSTAKAAN membahas penelitian terdahulu serta teori objek, fitur, model, dan evaluasi penelitian.
    \item BAB 3 METODOLOGI menjelaskan protokol pengumpulan dan audit data, target, konfigurasi fitur, pelatihan, serta evaluasi kronologis.
    \item BAB 4 HASIL menyajikan hasil audit, sampel evaluasi, kinerja model, perubahan berpasangan, dan analisis sensitivitas.
    \item BAB 5 PEMBAHASAN menghubungkan hasil dengan rumusan masalah, konsistensi kontribusi kelompok fitur, penelitian terdahulu, dan ancaman validitas.
    \item BAB 6 PENUTUP memuat kesimpulan serta saran berdasarkan hasil dan keterbatasan penelitian.
\end{enumerate}
